\documentclass[12pt,a4paper]{report}
% Общая преамбула. Подключается из subject/main.tex после \documentclass.

% ============================================================
% Базовые шрифты и кодировка
% ============================================================
\usepackage[T2A]{fontenc}
\usepackage[utf8]{inputenc}
\usepackage[russian]{babel}

\renewcommand{\rmdefault}{cmr}
\renewcommand{\sfdefault}{cmss}
\renewcommand{\ttdefault}{cmtt}
\renewcommand{\familydefault}{\rmdefault}

% ============================================================
% Математика и базовое оформление
% ============================================================
\usepackage{amsmath,amssymb,amsthm,mathtools}
\usepackage{mathrsfs}
\usepackage{bbold}
\usepackage{enumitem}
\usepackage{xparse}
\usepackage{xcolor}
\usepackage{microtype}
\usepackage{geometry}          % для настройки полей
\usepackage{parskip}           % для расстояний между абзацами (без параметров)
\usepackage{hyperref}
\usepackage[nameinlink,noabbrev]{cleveref}
\usepackage{tikz}
\usetikzlibrary{
    positioning,
    arrows.meta,
    calc,
    intersections,
    patterns,
    shapes.geometric,
    decorations.markings,
    matrix,
    fit,
    backgrounds
}

% ============================================================
% Рисунки, графики, схемы, таблицы и фрагменты кода
% ============================================================
\usepackage{graphicx}      % изображения
\usepackage{float}         % точное размещение рисунков/таблиц через [H]
\usepackage{caption}       % подписи к рисункам и таблицам
\usepackage{subcaption}    % несколько рисунков внутри одной фигуры
\usepackage{wrapfig}       % обтекание рисунков текстом
\usepackage{pgfplots}      % графики функций и численных данных
\pgfplotsset{compat=1.18}
\usepackage{tikz-cd}       % коммутативные диаграммы

\usepackage{booktabs}      % аккуратные таблицы
\usepackage{array}
\usepackage{tabularx}
\usepackage{longtable}
\usepackage{multirow}

\usepackage{listings}      % код: C/C++/Python/SQL и др.

% Дополнительная математика
\usepackage{bm}            % жирные математические символы: \bm{x}
\usepackage{bbm}           % красивый blackboard bold для цифр: \mathbbm{0}, \mathbbm{1}
\usepackage{esint}         % \oiint, \varoiint и другие интегралы
\usepackage{cancel}        % \cancel{}, \bcancel{}, \xcancel{}

% ============================================================
% Отступы и расстояния
% ============================================================
\geometry{a4paper,top=2cm,bottom=2cm,left=3cm,right=3cm,marginparwidth=1.75cm}
% parskip уже загружен без параметров, что даёт нулевой отступ первой строки
% и положительный межстрочный интервал между абзацами.

\setlength{\emergencystretch}{3em}

\hypersetup{
    colorlinks=true,
    linkcolor=black,
    urlcolor=black,
    citecolor=black,
    linktoc=all
}

% ============================================================
% Доказательства
% ============================================================
\addto\captionsrussian{%
    \renewcommand{\proofname}{Доказательство}%
}

\makeatletter
\renewenvironment{proof}[1][\proofname]{%
    \par
    \begingroup
    \leftskip=\dimexpr-\@totalleftmargin\relax
    \rightskip=0pt
    \parindent=0pt
    \noindent\textit{\underline{#1.}}\par\nobreak
    \noindent\ignorespaces
}{%
    \unskip\nobreak\hspace{0.15em}\ensuremath{\blacksquare}%
    \par
    \endgroup
}
\makeatother

% ============================================================
% Заголовки
% ============================================================
\usepackage{titlesec}
\titleclass{\chapter}{straight}

\titleformat{\chapter}[display]
  {\normalfont\huge\bfseries\raggedright}
  {\chaptername\ \thechapter}
  {0.7em}
  {\Huge\raggedright}

\titleformat{\section}
  {\normalfont\Large\bfseries}
  {\thesection}{0.7em}{}

\titleformat{\subsection}
  {\normalfont\large\bfseries}
  {\thesubsection}{0.7em}{}

\titleformat{\subsubsection}
  {\normalfont\normalsize\bfseries}
  {\thesubsubsection}{0.7em}{}

\titlespacing*{\chapter}{0pt}{0.45em}{0.45em}
\titlespacing*{\section}{0pt}{0.55em}{0.25em}
\titlespacing*{\subsection}{0pt}{0.45em}{0.18em}
\titlespacing*{\subsubsection}{0pt}{0.35em}{0.15em}

% ============================================================
% Лекции
% ============================================================
\newcounter{lection}
\NewDocumentCommand{\newlection}{m}{%
    \refstepcounter{lection}%
    \par\smallskip
    {\centering\color{gray!70}\normalsize Лекция \thelection. #1\par}
    \vspace{0.1em}
}

% ============================================================
% Теоремы, определения, замечания
% ============================================================
\makeatletter
\NewDocumentCommand{\definition}{o +m}{%
    \par
    \noindent
    \ifdim\@totalleftmargin>0pt
        \hspace*{\dimexpr-\@totalleftmargin\relax}%
    \fi
    \textbf{Определение\IfValueT{#1}{ (#1)}.}~#2
    \par
}
\makeatother

\ExplSyntaxOn

\NewDocumentCommand{\theorem}{o +m}{%
    \par\noindent
    \textbf{%
        Теорема
        \IfValueT{#1}{%
            \regex_match:nnTF { \A \d+ \Z } { #1 }
                {~#1}
                {~(#1)}
        }%
        .
    }~#2
    \par
}

\ExplSyntaxOff

\NewDocumentCommand{\proposal}{o +m}{%
    \par\noindent
    \textbf{Предложение\IfValueT{#1}{~#1}.}~#2
    \par
}

\NewDocumentCommand{\fact}{o +m}{%
    \par\noindent
    \textit{Факт\IfValueT{#1}{ (#1)}.}~#2
    \par
}

\NewDocumentCommand{\question}{+m}{%
    \par\noindent
    \textbf{Вопрос.}~#1
    \par
}

\NewDocumentCommand{\note}{+m}{%
    \par\noindent
    \textbf{Замечание.}~#1
    \par
}

\NewDocumentCommand{\corollary}{o +m}{%
    \par\noindent
    \textbf{Следствие\IfValueT{#1}{~#1}.}~#2
    \par
}

\NewDocumentCommand{\example}{+m}{%
    \par\noindent
    \textbf{Пример.}~#1
    \par
}

% ============================================================
% Списки
% ============================================================
\NewDocumentCommand{\numbers}{+m}{%
    \begin{enumerate}[leftmargin=2em,itemsep=0.25em,topsep=0.3em]
    #1
    \end{enumerate}
}

\NewDocumentCommand{\bullets}{+m}{%
    \begin{itemize}[leftmargin=2em,itemsep=0.25em,topsep=0.3em]
    #1
    \end{itemize}
}

\NewDocumentCommand{\provehere}{+m}{%
    \par\smallskip
    \begin{proof}#1\end{proof}
}

\NewDocumentCommand{\provewthen}{m +m}{%
    \par\smallskip
    \begin{proof}[#1]#2\end{proof}
}

\NewDocumentCommand{\provebullets}{+m}{%
    \begin{proof}
    \begin{itemize}[leftmargin=2em,itemsep=0.25em,topsep=0.3em]
    #1
    \end{itemize}
    \end{proof}
}

\NewDocumentCommand{\provenumbers}{+m}{%
    \begin{proof}
    \begin{enumerate}[leftmargin=2em,itemsep=0.25em,topsep=0.3em]
    #1
    \end{enumerate}
    \end{proof}
}

% ============================================================
% Дополнительные команды
% ============================================================
\usepackage{empheq}      % для \encircle
\usepackage{changepage}  % для \blockindent

% Леммы
\newtheorem{lemma_env}{Лемма}[section]
\newcommand{\lemma}[2][]{\begin{lemma_env}[#1]#2\end{lemma_env}}

% Отступ (переименовано из \indent, чтобы не ломать ядро LaTeX)
\newcommand{\blockindent}[1]{%
  \begin{adjustwidth}{1.5cm}{}#1\end{adjustwidth}%
}

% Одностраничные блоки
\newcommand{\singlepage}[1]{\begin{samepage}#1\end{samepage}}

% Сборки формул
\renewcommand{\gather}[1]{\begin{gather*}#1\end{gather*}}
\renewcommand{\multline}[1]{\begin{multline*}#1\end{multline*}}

% Обрамлённые формулы
\newcommand{\encircle}[1]{\begin{empheq}[box=\fbox]{align*}#1\end{empheq}}

% Математические операторы и символы
\newcommand{\proddots}{\cdot\ldots\cdot}
\newcommand{\abs}[1]{\left|#1\right|}
\newcommand{\angles}[1]{\left\langle#1\right\rangle}
\newcommand{\der}[2]{\frac{\partial #1}{\partial #2}}
\DeclareMathOperator{\differential}{d}
\newcommand{\dif}{\differential\!}
\newcommand{\fmax}{\lor}
\newcommand{\fmin}{\land}
\DeclareMathOperator*{\esssup}{ess\,sup}
\DeclareMathOperator*{\essssup}{ess\,sup}
\DeclareMathOperator*{\esssssup}{ess\,sup}
\newcommand{\Map}{\longrightarrow}
\newcommand{\almost}[1]{\underset{#1}{\overset{\textup{почти всюду}}\Map}}
\newcommand{\circlesign}[1]{\mathrel{\tikz[baseline=-0.5ex]{\node[draw, circle, inner sep=1pt] (char) {\vphantom{$-$}$#1$};}}}

% Дополнительные команды для доказательств
\newcommand{\indentlemma}[3][]{\blockindent{\lemma[#1]{#2\prove[Доказательство леммы]{#3}}}}
\newcommand{\prove}[2][Доказательство]{\begin{proof}[#1]#2\end{proof}}

% ============================================================
% Часто используемые математические обозначения
% ============================================================
\newcommand{\N}{\mathbb{N}}
\newcommand{\Z}{\mathbb{Z}}
\newcommand{\Q}{\mathbb{Q}}
\newcommand{\R}{\mathbb{R}}
\renewcommand{\C}{\mathbb{C}}
\newcommand{\Ord}{\mathrm{Ord}}
\newcommand{\Card}{\mathrm{Card}}

\renewcommand{\o}{\varnothing}
\newcommand{\bs}{\setminus}
\newcommand{\inv}{^{\complement}}
\newcommand{\then}{\Rightarrow}
\newcommand{\map}{\to}
\newcommand{\bydef}{\coloneqq}
\newcommand{\all}[1]{\left\{\begin{aligned}#1\end{aligned}\right.}
\newcommand{\any}[1]{\left[\begin{aligned}#1\end{aligned}\right.}
\newcommand{\exist}{\exists}
\newcommand{\dom}{\operatorname{dom}}
\newcommand{\rng}{\operatorname{rng}}
\newcommand{\defset}[2]{\left\{#1\;\middle|\;#2\right\}}
\newcommand{\switch}[1]{\left\{\begin{aligned}#1\end{aligned}\right.}

% ------------------------------------------------------------
% Большие операторы: индексы автоматически ставятся сверху/снизу
% ------------------------------------------------------------
\let\oldlim\lim
\renewcommand{\lim}{\oldlim\limits}
\let\oldsup\sup
\renewcommand{\sup}{\oldsup\limits}
\let\oldinf\inf
\renewcommand{\inf}{\oldinf\limits}
\let\oldmax\max
\renewcommand{\max}{\oldmax\limits}
\let\oldmin\min
\renewcommand{\min}{\oldmin\limits}

\let\oldsum\sum
\renewcommand{\sum}{\oldsum\limits}
\let\oldprod\prod
\renewcommand{\prod}{\oldprod\limits}
\let\oldcoprod\coprod
\renewcommand{\coprod}{\oldcoprod\limits}

\let\oldbigcap\bigcap
\renewcommand{\bigcap}{\oldbigcap\limits}
\let\oldbigcup\bigcup
\renewcommand{\bigcup}{\oldbigcup\limits}
\let\oldbigsqcup\bigsqcup
\renewcommand{\bigsqcup}{\oldbigsqcup\limits}

\let\oldbigvee\bigvee
\renewcommand{\bigvee}{\oldbigvee\limits}
\let\oldbigwedge\bigwedge
\renewcommand{\bigwedge}{\oldbigwedge\limits}
\let\oldbigoplus\bigoplus
\renewcommand{\bigoplus}{\oldbigoplus\limits}
\let\oldbigotimes\bigotimes
\renewcommand{\bigotimes}{\oldbigotimes\limits}
\let\oldbigodot\bigodot
\renewcommand{\bigodot}{\oldbigodot\limits}
\let\oldbiguplus\biguplus
\renewcommand{\biguplus}{\oldbiguplus\limits}

% Красивые blackboard-bold цифры (\mathbb не во всех шрифтах содержит цифры)
\newcommand{\bbzero}{\mathbbm{0}}
\newcommand{\bbone}{\mathbbm{1}}

% Универсальные обозначения для разных курсов
\newcommand{\norm}[1]{\left\lVert #1 \right\rVert}
\newcommand{\floor}[1]{\left\lfloor #1 \right\rfloor}
\newcommand{\ceil}[1]{\left\lceil #1 \right\rceil}
\newcommand{\paren}[1]{\left( #1 \right)}
\newcommand{\brackets}[1]{\left[ #1 \right]}
\newcommand{\set}[1]{\left\{ #1 \right\}}
\newcommand{\conj}[1]{\overline{#1}}
\newcommand{\vect}[1]{\bm{#1}}

\DeclareMathOperator{\RePart}{Re}
\DeclareMathOperator{\ImPart}{Im}
\DeclareMathOperator{\Arg}{Arg}
\DeclareMathOperator{\Ker}{Ker}
\DeclareMathOperator{\Image}{Im}
\DeclareMathOperator{\rank}{rank}
\DeclareMathOperator{\tr}{tr}
\DeclareMathOperator{\sgn}{sgn}
\DeclareMathOperator{\lcm}{lcm}
\DeclareMathOperator{\Span}{span}

% Differential notation in math; preserve the text dot-below accent.
\let\textdotbelow\d
\DeclareRobustCommand{\d}{\ifmmode\mathrm{d}\else\expandafter\textdotbelow\fi}

\endinput

% Всё ниже оставлено из исходного файла как справочный образец титульной
% страницы и LaTeX не читает после команды \endinput.
% ============================================================
% Информация о документе и титульная страница
% ============================================================
\date{III семестр, 2026 г.}
\title{Название предмета. Неофициальный конспект}
\author{Лектор: ФИО лектора\\
Конспект: Викторова Дарья, 25.Б03-пу}

\begin{document}

\begin{titlepage}
    \thispagestyle{empty}
    \begin{center}
        \vspace*{0.6cm}

        {\large Санкт-Петербургский государственный университет\par}
        \vspace{0.35em}
        {\normalsize Факультет прикладной математики~--- процессов управления\par}

        \vspace*{5.1cm}

        {\fontsize{30}{35}\selectfont\bfseries Название предмета\par}
        \vspace{0.65em}
        {\fontsize{17}{21}\selectfont\itshape Неофициальный конспект\par}

        \vspace{2.7cm}

        {\normalsize
        \begin{tabular}{@{}rl@{}}
            \textit{Лектор:} & ФИО лектора \\
            \textit{Конспект:} & Викторова Дарья, 25.Б03-пу
        \end{tabular}\par}

        \vfill

        {\large III семестр\par}
        \vspace{0.22em}
        {\large 2026 год\par}
        \vspace*{0.55cm}
    \end{center}
\end{titlepage}

\tableofcontents
\newpage

\setcounter{lection}{0}
\newlection{Число Месяц 2026 г.}

\end{document}



\title{Теория групп и теория чисел. Неофициальный конспект}
\author{Лектор: Кривошеин Александр Владимирович\\
Конспект: Викторова Дарья, 25.Б03-пу}
\date{III семестр, 2026 г.}

\begin{document}

\begin{titlepage}
    \thispagestyle{empty}
    \begin{center}
        \vspace*{0.6cm}

        {\large Санкт-Петербургский государственный университет\par}
        \vspace{0.35em}
        {\normalsize Факультет прикладной математики~--- процессов управления\par}

        \vspace*{5.1cm}

        {\fontsize{30}{35}\selectfont\bfseries Теория групп и теория чисел\par}
        \vspace{0.65em}
        {\fontsize{17}{21}\selectfont\itshape Неофициальный конспект\par}

        \vspace{2.7cm}

        {\normalsize
        \begin{tabular}{@{}rl@{}}
            \textit{Лектор:} & ФИО лектора \\
            \textit{Конспект:} & Викторова Дарья, 25.Б03-пу
        \end{tabular}\par}

        \vfill

        {\large III семестр\par}
        \vspace{0.22em}
        {\large 2026 год\par}
        \vspace*{0.55cm}
    \end{center}
\end{titlepage}

\tableofcontents
\newpage

\setcounter{lection}{0}

\newlection{3 сентября 2026 г.}

\chapter{Теория групп}

\textit{Теория групп} - описание симметрии.

\definition{\textit{Симметрия} - биективные преобразования над множеством объектов, которые сохраняют какие-то свойства.}

Пример. 4 штуки симметрий (поворот на $180^0$, отражение, параллельный перенос, тождественное преобразование), сохраняющих расстояние между любыми двумя точками (изометрия).

Добавить рисунок: ДПСК и нарисован квадрат на ДПСК

\section{Предварительные сведения}

Отображение. $f\colon X\to Y$.
$D\subset X, f(D)$
$E\subset Y, f^{-1}(E)$

Композиция. $f\colon X\to Y$.
$g\colon Y\to Z$.
$g \circ f\colon X\to Z$
$(g\circ f)(x)=g(f(x))$

Ассоциативность композиции.
\[h\circ(g\circ f)=(h\circ g)\circ f\]

Инъекция
Добавить рисунок: инъекция

Сюръекция
Добавить рисунок: сюръекция

Биекция = инъекция + сюръекция.

$f\colon X\to Y$
Уравнение относительно $x\colon y = f(x)$.
\begin{enumerate}
    \item f - инъективна: если $\exists$, то !;
    \item f - сюръективна: всегда $\exists$, но может быть не !;
    \item f - биективна: $\exists!$ решение.
\end{enumerate}

$id_x\colon X\to X$ - тождественное отображение, $id_x(x)=x$.

\definition{
Отображение $f\colon X\to X$ \textit{обратимо}, если $\exists g\colon Y\to X\colon g\circ f=id_x, f\circ g = id_y$.

g называется \textit{обратным отображением}.
}

Мини-теорема. Обратное отображение единственно.
\begin{proof}
    Пусть существует два отображения $g_1,g_2$. \\
    $g_1=g_1\circ id_y=g_1\circ(f\circ g_2)=(g_1\circ f_\circ g_2=id_x\circ g_2=g_2$.
\end{proof}

Утверждение. $f$ обратимо $\iff$ $f$ биективно.
Обозначим $f^{-1}$ - обратное к $f$.

\definition{$X$. Бинарное отношение $R$ на $X$ - это $R\subset X\times X$.}

Пример. Бинарное отношение победы первого игрока в игре "Камень, ножницы, бумага".

Добавить рисунок

\definition{Будем называть R \textit{отношением эквивалентности}, если оно
\begin{enumerate}
    \item реффлексивно: $XRX$;
    \item симметрично: $XRY\Rightarrow YRX$;
    \item транзитивно: $XRY, YRZ\Rightarrow XRZ$
\end{enumerate}

Обозначение: XRY.
Запись отношения эквивалентности: $X\sim Y$.}

Классы эквивалентности.
\[[x]=\{y\in X\colon x\sim y\}\]
$X=\bigcup_{x - \text{представители классов}} [x]$

Классы эквивалентности либо совпадают, либо не пересекаются.

Пример. Классы вычетов.
\[ Z/mZ=\{[k], k = \overline{0, m-1}\}, \text{где} [k]=k+mZ. \]

\section{Определение группы. Базовые свойства.}

\textit{Бинарная операция} на $X$ - $\tau \colon X\times X\to X$, \\
$(X,\tau)$ - \textit{алгебраическая структура} (magma).

\definition{Группа G - это алгебраическая структура вида $(G, *)$, где бинарная операция удовлетворяет свойствам:
\begin{enumerate}
    \item Ассоциативности: $a *(b*c)=(a*b)*c \forall a,b,c \in G$;
    \item Существование нейтрального элемента $e\colon a * e = e* a = a\forall a\in G$.
    \item $\forall a \in G\exists b \in G\colon a * b = b * a = e$.
\end{enumerate}

Группа называется \textit{абелевой}, если бинарная операция коммутативна:
\[a * b = b * a\forall a,b\in G\]
}

Мультипликативная запись: $*\to \cdot$;
Аддитивная запись: $*\to +$.

\subsection{Свойства группы}

\begin{enumerate}
    \item G - непуста;
    Минимальная по размеру группа: $G=\{e\}$.
    \item Нейтральный элемент единственен.
    \item Обратный элемент единственен.
    Для $a\in G$ обозначим обратный элемент: $a^{-1}$.
    \item $(a^{-1})^{-1}=a$.
    \begin{proof}
        К $a^{-1}$ обратным является $a$; к $a^{-1}$ обратным является $(a^{-1})^{-1}$. Следовательно, $(a^{-1})^{-1}=a$.
    \end{proof}
    \item $(a\cdot b)^{-1} = b^{-1}\cdot a^{-1}$.
    \item Гарантируется решение уравнений вида $ax=b$;
    Решение: $x=a^{-1}\cdot b$.
\end{enumerate}

\definition{
Порядок группы G - количество элементов в ней, если она конечна; $+\infty$, если группа содержит бесконечное число элементов.

Обозначение: $|G|$.
}

\subsection{Таблицы умножения конечных групп (таблицы Кэли)}
Перенумеруем элементы $G=\{e, g_2,\dots,g_n\}$, $|G|=n$.

добавить рисунок таблицы

\theorem{Таблицы Кэли являются латинским квадратом, то есть в любой строке, в любом столбце таблицы все элементы различны.}
\begin{proof}
    Предположим, что существует два одинаковых элемента.
    \[\begin{pmatrix}g_k\cdot e & g_k\cdot g_2 & \dots & g_k\cdot g_k & \dots &g_k\cdot g_n\end{pmatrix}.\]
    $g_k\cdot g_i=g_k\cdot g_j, i\neq j$.
    Домножим обе части на $g_k^{-1}$.
    Получим, что $g_i=g_j$. Противоречие.
\end{proof}

капец написать))

\definition{
Пусть $(G, \cdot), (H, *)$ --- группы. \\
Говорят, что $G$ \textit{изоморфно} $H$, если $\exists \varphi\colon G\to H$ --- биекция. \\
\textit{Обозначение}: G $\cong$ H. \\
Для $a, b \in G$:
\[
\underbrace{\varphi(a\cdot b) = \varphi(a) * \varphi(b)}_{\text{$\varphi$ сохраняет операцию}} \quad (1)
\]
}

\[
\begin{array}{c|cccc}
 & 0 & 1 & 2 \\ \hline
0 & 0 & 1 & 2 \\
1 & 1 & 2 & 0 \\
2 & 2 & 0 & 1
\end{array}
\qquad
\qquad
\qquad
\begin{array}{c|cccc}
 & e & a & b \\ \hline
e & e & a & b \\
a & a & b & e \\
b & b & e & a
\end{array}
\]

\newlection{4 сентября 2026 г.}

\subsection{Свойства изоморфизма}
Пусть $G, H, K$ --- группы, $G\cong H$. Тогда
\begin{enumerate}
    \item $|G|=|H|$.
    \item $\varphi$ --- изоморфизм. Тогда $\varphi^{-1}$ --- изоморфизм, и $H \cong G$.
    \begin{proof}
        Проверим равенство $(1)$ для $\varphi^{-1}$, $\varphi^{-1}\colon H\to G$. \\
        Рассмотрим $U, V \in H$. Надо доказать, что 
        $\varphi^{-1}(U\cdot V)=\varphi^{-1}(U)\cdot\varphi^{-1}(V)$. \\
        Существуют $a, b\in G$ такие, что $\varphi(a)=U$, $\varphi(b)=V$. 
        Тогда
        \[
        \varphi^{-1}(U\cdot V)=\varphi^{-1}(\varphi(a)\cdot \varphi(b))=\varphi^{-1}(\varphi(a\cdot b))=a\cdot b=\varphi^{-1}(a)\cdot \varphi^{-1}(b).
        \]
    \end{proof}
    \item $\varphi\colon G\to H, \psi\colon H \to K$ --- изоморфизмы. Тогда $\varphi \circ\psi$ --- изоморфизм.
    \[
    \begin{cases}
        G \cong H \\ H \cong K
    \end{cases}
    \Rightarrow G\cong K.
    \]
    \begin{proof}
        $a, b \in G$.
        \[
        (\underbrace{\varphi \circ \psi}_h)(a\cdot b)=\varphi(\psi(a\cdot b))=\varphi(\psi(a)\cdot \psi(b))=\varphi(\psi(a))\cdot \varphi(\psi(b))=(\underbrace{\varphi \circ \psi}_h)(a)\cdot (\underbrace{\varphi \circ \psi}_h)(b).
        \]
    \end{proof}
    \item $G \cong G$.
    \item $\varphi(e_G)=e_H$.
    \begin{proof}
        $a\in G$.
        \[
        \varphi(a)=\varphi(a\cdot e_G)=\varphi(a)\cdot \varphi(e_G).
        \]
        Домножим слева на $(\varphi(a))^{-1}$. Получим $e_H=\varphi(e_G)$.
    \end{proof}
\end{enumerate}

\example{%
$(\mathbb R, +)\cong(\mathbb R_{\underbrace{+}_{>0}}, \cdot)$. \\

\begin{center}
\begin{tikzpicture}[scale=0.6]
    \draw[->] (-3,0) -- (3,0) node[below]{$x$};
    \draw[->] (0,-1) -- (0,5) node[left]{$y$};
    \draw[red, thick, domain=-2.5:2, smooth, samples=50] plot(\x, {exp(\x)});
\end{tikzpicture}
\end{center}

\[
e^{x+y}=e^x\cdot e^y.
\]
}

\example{
$\mathbb V$ --- конечномерное векторное линейное пространство над $\mathbb R$. \\
$A\colon \mathbb V \to \mathbb V$, $A$ --- невырожденный линейный оператор.
\[
A(\alpha x+\beta y)=\alpha Ax+\beta Ay.
\]
$GL(\mathbb V)=\{A\colon \mathbb V\to\mathbb V \mid A \text{ --- невырожденный линейный оператор}\}$ --- группа относительно композиции.
}

\subsection{Операция композиции}
\[
(A\circ B)x=ABx.
\]
\begin{enumerate}
\item $Ix=x$ --- нейтральный элемент;
\item $(GL(\mathbb V), \circ)$ --- группа;
\item $GL(\mathbb V)\cong GL_{\dim \mathbb V}(\mathbb R)$;
\item $\varphi(A)$ --- матрица оператора в выбранном базисе.
\end{enumerate}

\note{
Набор симметрий некоторого объекта образует группу.
}

\subsection{Группа симметрий}
\[
\operatorname{Sym}(X)=\{f\colon X\to X \mid f \text{ --- биекция, сохраняющая структуру}\}.
\]
\begin{enumerate}
\item $(f\circ g)(x)=f(g(x))$ --- композиция симметрий;
\item $e=\mathrm{id}_X$ --- нейтральный элемент;
\item $(\operatorname{Sym}(X), \circ)$ --- группа.
\end{enumerate}

пример
пример

\section{Подгруппы}

\definition{
Пусть $G$ - группа, $H\subset G$ называется \textit{подгруппой}, если $H$ - группа с той же бинарное операцией, что и в $G$. \\
\textit{Обозначение}: $H\leq G$ или $H<G$, если $\begin{cases}
    H \leq G \\
    H \neq G
\end{cases}$
.
}

\theorem[Критерий "быть подгруппой"]{
\[
H\leq G \iff
\]
\begin{enumerate}
\setcounter{enumi}{-1}
\item $H\neq\varnothing$;
\item Замкнутость по умножению: $\forall a,b\in H\Rightarrow a\cdot b\in H$;
\item Замкнутость по взятию обратного элемента: $\forall a\in H\Rightarrow a^{-1}\in H$.
\end{enumerate}
}

\begin{proof}
    \begin{enumerate}
        \item[$\Longrightarrow$] \textit{Необходимость}. 
        Если $H\leq G$, то $H$ — группа. Следовательно, все три условия выполняются по определению группы.
        
        \item[$\Longleftarrow$] \textit{Достаточность}. 
        \begin{enumerate}
            \item В $H$ есть нейтральный элемент. \[
            \forall a\in H \exists a^{-1}\in H \colon a\cdot a^{-1}=e\in H. 
            \]
            
            \item Ассоциативность операции наследуется. $\forall a,b,c \in H\subseteq G\colon$
            \[
            a\cdot (b\cdot c)=(a\cdot b)\cdot c.
            \]
            \item Обратный элемент $\forall a\in H$ существует по условию (2).
        \end{enumerate}
        Следовательно, $H$ — группа, значит, $H\leq G$.
    \end{enumerate}
\end{proof}

\fact{В конечной группе условие замкнутости по умножению автоматически влечёт замкнутость по взятию обратного. Поэтому для конечного подмножества достаточно проверить только непустоту и умножение.}

\note{
Для любой группы $G$ существуют тривиальные подгруппы:
\[
G\leq G \quad\text{и}\quad \{e\}\leq G.
\]
}

\example{}
\example{}

\definition{
Рассмотрим $G$, $M\subset G$.
\textit{Подгруппой, наименьшей по включению и содержащей $M$}, называется $\langle M \rangle=\bigcap\limits_{\substack{H\leq G \\ M \subset H}}H.$ \\

Пересечение подгрупп --- подгруппа. \\

Введем обозначения:
\begin{enumerate}
    \item $g^k=g^{k-1}\cdot g$, $k\in\mathbb N$, где $g^1=g$.
    \item $g^0 = e$.
    \item $g^{-k} := (g^k)^{-1}$.
\end{enumerate}
}

\theorem[характеризация $\langle M\rangle$]{
\[
\langle M\rangle = \underbrace{\left\{ e\cdot g_1^{\varepsilon_1}\cdot g_2^{\varepsilon_2}\cdots g_n^{\varepsilon_n} \;\middle|\; n\in\mathbb N,\; g_i\in M,\; \varepsilon_i\in\{\pm1\} \right\}}_V.
\]
(\textit{Повторения $g$ разрешаются}).
}
\begin{proof}
    \begin{enumerate}
        \item Надо доказать, что $\langle M\rangle \subset V$. \\
        Ясно, что $\forall g_i\in M\Rightarrow g_i\in V$. Докажем, что $V$ --- подгруппа. Возьмем $a,  b\in V\Rightarrow a\cdot b\in V$. Рассмотрим
        \[
        (eg_1^{\varepsilon_1}\cdot\dots\cdot g_n^{\varepsilon_n})^{-1}=g_n^{-\varepsilon_n}\cdot \dots \cdot g_1^{-\varepsilon_1}\cdot e^{-1}=e\cdot g_n^{-\varepsilon_n}\cdot \dots \cdot g_1^{-\varepsilon_1}\in V.
        \]
        Получили, что $V$ --- подгруппа в $G$ и $M\subset V$. Следовательно, $\langle M\rangle \subset V$.
        \item  Надо доказать, что $V\subset \langle M\rangle$. \\
        $\langle M\rangle$ --- подгруппа. Следовательно,  $\langle M\rangle$ содержит в себе всевозможные произведения своих элементов и их обратных. Откуда $V\subset \langle M\rangle$.
    \end{enumerate}
    Следовательно, $V=\langle M\rangle$.
\end{proof}

Пусть $G$ --- группа, $g\in G$.
\definition{
Подгруппа $\langle g\rangle$ называется \textit{циклической подгруппой} группы $G$.
}

\definition{
Группа $G$ называется \textit{циклической}, если $\exists g\in G\colon \langle g\rangle = G$.
\[
\langle g\rangle = \{g^n \mid n\in \mathbb Z\}.
\]
}

\example{}

\section{Группа остатков}

Рассмотрим множество $\mathbb{Z}_m\setminus\{0\}$ с операцией умножения по модулю $m$:
\[
a *_m b := a\cdot b \pmod m.
\]

Возникает вопрос: является ли $(\mathbb{Z}_m\setminus\{0\}, *)$ группой?

\example{}

\theorem{
$Z_m\setminus\{0\}$, с операцией $*_m$.
\[
a \text{ --- имеет обратный элемент по умножению} \iff \text{НОД}(a,m)=1.
\]
}
\begin{proof}
    \begin{enumerate}
        \item[$\Longleftarrow$] \textit{Достаточность}. 
        Пусть $\text{НОД}(a,m)=1$. Рассмотрим линейное представление:
        \[
        U\cdot a+V\cdot m = 1.
        \]
        Получим, что $U\cdot a \equiv 1 \pmod{m} \Rightarrow U=qm+u\pmod m\Rightarrow (U\bmod m)\cdot a\equiv 1\pmod m$.\\ 
        Следовательно, $U\bmod m$ --- обратный по умножению к $a$.

        \item[$\Longrightarrow$] \textit{Необходимость}. 
        Пусть $b$ --- обратный по умножению к $a$. Тогда
        \[
        a *_m b = \underbrace{ab\pmod m}_{ab=qm+1}=1,
        \]
        \[
        \underbrace{b}_U a - \underbrace{q}_V m = 1.
        \]
        Следовательно, $\text{НОД}(a,m)=1$.
    \end{enumerate}
\end{proof}

\end{document}


